\section{Creación de Geometría}
La creación de la geometría se realiza utilizando grupos de nodos de geometría, creados a través
del editor de nodos de geometría. Cada grupo contiene parametros expuestos que controlan
proceduralmente la creación de la geometría, y por lo tanto, son valores que pueden ser
animados en el editor de Blender. Por consecuencia, los grupos actúan también como
controles de animación.

\subsection{Visualización de Volumenes}
La visualización de volúmenes se realiza a través de un grupo de nodos \textit{Volume Viewer}.
Siendo que los tensores son codificados en la forma de volúmenes de OpenVDB, su visualización consiste
principalmente en el muestreo de los valores en la cuadrícula tridimensional.
\medskip

El grupo de nodos toma como parámetros:
\begin{itemize}
    \item \textit{Volume File}: El archivo OpenVDB con el tensor capturado de una capa de la red.
    \item \textit{Resolution}: La resolución de una rebanada del tensor. Se asume que las
    rebanadas son cuadradas.
    \item \textit{Slice Count}: El número de rebanadas del tensor, o el número de canales que produce.
    \item \textit{World Size}: El tamaño físico que la visualización tendrá en la escena de Blender. Las dimensiones
    X, Y y Z corresponden a la anchura, altura y grosor de la visualización del volumen dentro de la escena.
    \item \textit{Slice Selection}: El índice de la rebanada a visualizar. La visualización oculta todas
    las rebanadas frente a la selección.
    \item \textit{Single Slice}: Opción para ocultar todas las rebanadas excepto la seleccionada.
\end{itemize}

El nodo de geometría genera instancias de planos que actúan como las visualizaciones de cada revanada. 
Las dimensiones de cada plano estan dadas por los valores especificados por \textit{World Size}, y subidivido 
de acuerdo a la resolución. Se intancia uno de estos planos por rebanada en el volumen y se le asigna un índice
correspondiente.
\medskip

Seguido, haciendo uso del mapa UV autogenerado por cada plano y su índice de rebanada se muestrea en
el volumen el voxel correspondiente. El valor obtenido, correspondiente a un valor de activacion, 
es pasado al material asignado a los planos, un valor por cada subdivisión, para mapear el valor a un color.
\medskip

Con lógica adicional el nodo de geometría genera una malla para visualizar los límites físicos del volumen
y selectivamente ocultar capas. Y finalmente se almacenan como atributos del objeto las dimensiones físicas
del volumen, las dimensiones del tensor (resolución y número de capas) y el índice de la capa seleccionada.

\subsection{Visualización de Conexiones}
La visualización de conexiones se implementa a través del grupo de nodos \textit{Volume Connection}, el cual
realiza dos tareas principales: visualiza físicamente las conexiones entre capas, y visualiza el flujo de 
datos entre estas.
\medskip

El grupo toma como parámetros:
\begin{itemize}
    \item \textit{Input Layer, Output Layer}: Referencias a los objetos de visualización de volúmenes
    corespondientes a capas.
    \item \textit{View Convolution}: Muestra la animación procedural de una convolución, mostrando un kernel
    iterando sobre las capas de entrada y salida.
    \item \textit{T}: Valor de progresión de la animación de la convolución.
    \item \textit{Kernel Size}: Tamaño del kernel.
    \item \textit{Stride}: Stride de la convolucion.
    \item \textit{Ceiling Mode}: Modo de redondeo cuando el tamaño del tensor es reducido debido al stride.
    Por defecto utiliza la operación \textit{ceiling}, siendo la otra opción \textit{floor}.
    \item \textit{Operation Mode}: Indica el modo de la operación de convolución, siendo los valores \textit{One to One},
    \textit{Many to One} y \textit{Many to Many}.
\end{itemize}

Para visualizar las conexiones físicas simplemente genera aristas entre las esquinas correspondientes de los 
volumenes, utilizando el atributo de las dimensiones físicas almacenadas por el nodo de geometría de visualización 
de volumenes de los objetos que representan las capas de entrada y salida.
\medskip

Para visualizar el flujo de datos los nodos generan una representación de la celda que la capa de entrada está 
evaluando y una representadión de la celda en la capa de salida con el valor de salida correspondiente. Utilizando
los atributos de las dimensiones físicas del volumen y las dimensiones del tensor se calcula el tamaño
físico de las celdas de una rebanada y la posición de la primera celda (en la esquina superior izquierda).
En el caso de la capa de entrada se genera también una representación de la localidad del kernel.
\medskip

Utilizando un valor \textit{T} se controla la progresión, de izquierda a derecha y de arriba a abajo, 
del kernel sobre las celdas de una rebanada, donde un valor de 0 indica la posición inicial y 1 la posición final 
(en la esquina inferior derecha). Utilizando los valores de la
resolución de la capa y el valor de stride se calcula la cuadrícula de iteración para representar los valores que
la capa evalúa, y utilizando el tamaño físico de las celdas se obtiene el desfase físico para trasladar correctamente 
el kernel. Cabe mencionar que las iteraciones de las capas de entrada y salida son realizadas independientemente,
y se asumie que los tamaños de tensores de entrada y salida son correctos para los parametros de \textit{stride} y 
\textit{ceil mode}.

Adicionalmente los kernels se pueden visualizar cubriendo todas las rebanadas (para operaciones que realizan 
operaciones sobre todas las capas, como la convolución), o solo la rebanada seleccionada 
(para capas que realizan operaciones por canal, como \textit{Max-Pooling}). A partir de estos se implementan
tres modos de visualización de las operaciones:
\begin{itemize}
    \item \textit{One to One}: Operaciones que se realizan canal a canal: cada canal produce un canal de salida,
    como en el caso de la operación de \textit{Max-Pooling}.
    \item \textit{Many to One}: La operación se aplica a todos los canales de entrada y produce un solo canal de
    salida, como en el caso de las capas convolucionales.
    \item \textit{Many to Many}: La intención de este modo es abstraer la visualización del flujo de datos entre
    capas, para simplemente visualizar la entrada y salida de las capas.
\end{itemize}

El kernel genera también un "fatasma", una copia del kernel posicionada directamente sobre la rebanada seleccionada
para auxiliar en la visualización de los valores que se evalúan.

\subsection{Visualización de Cajas}
El grupo de nodos \textit{Box Viewer} se utiliza para visualizar las cajas de salida de la red. Este grupo de nodos de 
geometría toma como parámetros:
\begin{itemize}
    \item \textit{Box File}: El archivo OpenVDB que contiene los datos codificados de las cajas.
    \item \textit{Box Count}: El número total de cajas de predicción producidas por la red.
    \item \textit{Class List}: La lista de etiquetas de clases usadas por la red, separadas por comas.
    \item \textit{Box Selection}: Control para visualizar cajas individuales o por parámetros.
    \item \textit{Confidence}: Filtro de visualización de cajas en base a su objetividad (\textit{Objectiveness}).
    \item \textit{Class Probability}: Filtro de visualización de cajas en base a su probabilidad de clase.
\end{itemize}
Para generar las cajas primero se genera una "libreria" de mallas de texto con las etiquetas de las clases
a partis de la lista de etiquetas de los parametros.
Se instancian cuadrados que representan cada una de las cajas, cada uno con un índice. Con este índice se indexa
del volúmen OpenVDB loa datos codificados de cada caja. Con ellos, se aplican transformaciones a los cuadros
para colocarlos en las coordenadas de predicción y para escalarlos a las dimensiones apropiadas. Utilizando 
el ID de clase de las cajas se insancia la malla de texto de la etiqueta de clase correspondiente sobre la
caja, y se asigna un color a la etiqueta y caja para facilitar la distinción de clases y se guarda como 
atributo la objetividad y probabilidad de clase por cada par de etiqueta/caja.
Finalmente se instancia un plano con la imagen de entrada sobre el cual se superponen las cajas y etiquetas
para visualizar la salida de la red.
\medskip

Para visualizar las caja hay tres opciones dependiendo del valor de \textit{Box Selection}:
\begin{itemize}
    \item \textit{-1}: Muestra todas las cajas, pero filtrando de acuerdo al valor de \textit{Confidence}
    y \textit{Class Probability}.
    \item \textit{0 - (N-1)}: Muestra únicamente la caja con el indice seleccionado, para una cuenta total de
    \textit{N} cajas.
    \item \textit{N}: Muestra todas las cajas, ignorando los valores de \textit{Confidence} y 
    \textit{Class Probability}.
\end{itemize}